\notechapter{N-20220412}{Indices, Dual Maps, Adjoints, Trace, and Singular Values}

\section{Indices carry type information before they abbreviate sums}

Let \(V\) have basis \(e_1,\ldots,e_n\), with dual basis \(e^1,\ldots,e^n\).  The Kronecker delta is
\[
  \delta_i^{\ j}
  =
  \begin{cases}
    1,&i=j,\\
    0,&i\ne j.
  \end{cases}
\]
It is the coordinate representation of the identity map:
\[
  I(e_j)=e_i\delta_i^{\ j}=e_j.
\]
The position of an index records which kind of slot it occupies.  A vector has components \(x^j\), a covector has components \(\alpha_i\), and contraction pairs one of each type:
\[
  \alpha(x)=\alpha_i x^i.
\]

Einstein convention suppresses the summation sign when an index occurs once upstairs and once downstairs:
\[
  \alpha_i x^i
  =\sum_i\alpha_i x^i.
\]
An index that survives once in every term is free and determines the output type.  A repeated index is dummy and may be renamed.  Expressions such as
\[
  A_i^{\ j}x^i
\]
are ill typed if \(i\) is intended to be the output index, because it occurs twice while \(j\) remains unmatched.  Index notation is valuable precisely because many algebraic mistakes become visible as type errors.

\section{Matrix multiplication is contraction of an output with an input}

Let
\[
  B:U\to V,
  \qquad
  A:V\to W.
\]
Choose bases and write
\[
  B(e_k)=B_j^{\ k}v_j,
  \qquad
  A(v_j)=A_i^{\ j}w_i.
\]
Then
\[
\begin{aligned}
  (A\circ B)(e_k)
  &=B_j^{\ k}A(v_j)\\
  &=A_i^{\ j}B_j^{\ k}w_i.
\end{aligned}
\]
Thus
\[
  (AB)_i^{\ k}=A_i^{\ j}B_j^{\ k}.
\]
The repeated index \(j\) is the intermediate vector-space direction that is produced by \(B\) and consumed by \(A\).

The identity tensor acts neutrally because
\[
  A_i^{\ j}\delta_j^{\ k}=A_i^{\ k},
  \qquad
  \delta_i^{\ j}A_j^{\ k}=A_i^{\ k}.
\]
Trace is another contraction:
\[
  \operatorname{tr}A=A_i^{\ i}.
\]
Matrix multiplication and trace are therefore not unrelated recipes; both pair an output index with a compatible input index.

\section{A change of basis transforms both indices and leaves contractions unchanged}

Let \(P\) be the change-of-basis matrix, so coordinate columns satisfy
\[
  [x]_{\mathrm{old}}=P[x]_{\mathrm{new}}.
\]
An endomorphism matrix transforms by similarity:
\[
  A'=P^{-1}AP.
\]
In indices,
\[
  A'_i{}^{j}
  =(P^{-1})_i{}^{r}
   A_r{}^{s}P_s{}^{j}.
\]
Contracting the free indices gives
\[
\begin{aligned}
  A'_i{}^{i}
  &=(P^{-1})_i{}^{r}
    A_r{}^{s}P_s{}^{i}\\
  &=A_r{}^{s}
    P_s{}^{i}(P^{-1})_i{}^{r}\\
  &=A_r{}^{s}\delta_s{}^{r}\\
  &=A_r{}^{r}.
\end{aligned}
\]
Hence
\[
  \operatorname{tr}(P^{-1}AP)
  =\operatorname{tr}A.
\]
The diagonal sum is invariant because it is a complete contraction: the basis-change matrix and its inverse cancel along the closed index loop.

By contrast, an individual diagonal entry \(A_i{}^i\) with \(i\) fixed is not invariant.  Only the sum over the repeated index has intrinsic meaning.

\section{Transpose is the matrix of the dual map}

Let
\[
  A:V\to W
\]
and define the dual map
\[
  A^\vee:W^*\to V^*,
  \qquad
  A^\vee(\varphi)=\varphi\circ A.
\]
Choose bases \(e_j\) of \(V\) and \(f_i\) of \(W\), with
\[
  A(e_j)=\sum_i A_{ij}f_i.
\]
For the dual basis vector \(f^i\),
\[
\begin{aligned}
  A^\vee(f^i)(e_j)
  &=f^i(Ae_j)\\
  &=A_{ij}.
\end{aligned}
\]
Therefore
\[
  A^\vee(f^i)=\sum_jA_{ij}e^j.
\]
The matrix of \(A^\vee\) has entry \(A_{ij}\) in row \(j\), column \(i\), so it is
\[
  [A^\vee]=A^T.
\]

The reversal law now follows without entry manipulation.  For
\[
  U\xrightarrow{B}V\xrightarrow{A}W,
\]
one has
\[
  (A\circ B)^\vee
  =B^\vee\circ A^\vee.
\]
In matrices,
\[
  (AB)^T=B^TA^T.
\]
Transposition reverses order because a functional is pulled backward through a composite map.

\section{An adjoint is a dual map returned to vectors by an inner product}

The transpose needs only vector spaces and duality.  An adjoint additionally needs inner products.  Let
\[
  A:V\to W
\]
be a linear map between finite-dimensional complex inner-product spaces.  The adjoint
\[
  A^\dagger:W\to V
\]
is characterized by
\[
  \langle Ax,y\rangle_W
  =\langle x,A^\dagger y\rangle_V.
\]
Riesz representation identifies each covector with a unique vector, so the adjoint is the composite
\[
  W\xrightarrow{\text{Riesz}}W^*
  \xrightarrow{A^\vee}V^*
  \xrightarrow{\text{Riesz}^{-1}}V.
\]

In orthonormal bases,
\[
  [A^\dagger]=\overline A^{\,T}=A^*.
\]
The operations should be kept distinct:
\[
\begin{array}{c|c}
\text{operation}&\text{matrix}\\ \hline
\text{transpose}&A^T\\
\text{entrywise conjugate}&\overline A\\
\text{conjugate transpose}&\overline A^{\,T}\\
\text{adjoint in orthonormal bases}&A^*.
\end{array}
\]
Over \(\mathbb R\), conjugation disappears and the adjoint matrix in orthonormal bases is the transpose.

\section{In nonorthonormal coordinates the Gram matrices enter the adjoint}

Suppose the chosen bases have Gram matrices
\[
  G_V=(\langle e_i,e_j\rangle_V),
  \qquad
  G_W=(\langle f_i,f_j\rangle_W).
\]
For coordinate columns,
\[
  \langle x,z\rangle_V=x^*G_Vz,
  \qquad
  \langle y,w\rangle_W=y^*G_Ww.
\]
The defining identity for the adjoint becomes
\[
  x^*A^*G_Wy
  =x^*G_V[A^\dagger]y
\]
for all \(x,y\).  Hence
\[
  \boxed{
  [A^\dagger]
  =G_V^{-1}A^*G_W.}
\]
The formula \(A^\dagger=A^*\) is therefore a feature of orthonormal coordinates, not the definition of adjointness.

For example, on \(\mathbb R^2\) use the inner product with
\[
  G=
  \begin{pmatrix}
    1&0\\
    0&4
  \end{pmatrix}
\]
and the linear map
\[
  A=
  \begin{pmatrix}
    0&1\\
    1&0
  \end{pmatrix}.
\]
Although \(A^T=A\), its adjoint in this metric is
\[
\begin{aligned}
  A^\dagger
  &=G^{-1}A^TG\\
  &=
  \begin{pmatrix}
    0&4\\
    1/4&0
  \end{pmatrix}.
\end{aligned}
\]
The coordinate swap is not self-adjoint when the second coordinate has four times the metric weight.

\section{Rank is the dimension of the image and the size of a minimal factorization}

For
\[
  A:V\to W,
\]
define
\[
  \operatorname{rank}A=\dim\operatorname{im}A.
\]
The first isomorphism theorem gives
\[
  V/\ker A\cong\operatorname{im}A.
\]
Thus rank is the dimension of the effective quotient after directions killed by \(A\) have been removed.

If \(\operatorname{rank}A=r\), choose a basis of the column space and place it in a matrix
\[
  B\in k^{m\times r}.
\]
Express every column of \(A\) in that basis and collect the coordinates in
\[
  C\in k^{r\times n}.
\]
Then
\[
  A=BC,
  \qquad
  \operatorname{rank}B
  =\operatorname{rank}C=r.
\]
Conversely, any factorization through an \(r\)-dimensional space has rank at most \(r\).  Hence rank is the smallest intermediate dimension through which the map can factor.

This factorization also proves equality of row and column rank.  Since \(A=BC\), every row of \(A\) is a linear combination of rows of \(C\), so
\[
  \operatorname{row}(A)
  \subseteq\operatorname{row}(C).
\]
Because \(B\) has full column rank, it has a left inverse \(L\) with
\[
  LB=I_r.
\]
Hence
\[
  C=LA,
\]
so every row of \(C\) is a linear combination of rows of \(A\).  Thus
\[
  \operatorname{row}(A)
  =\operatorname{row}(C),
\]
and its dimension is \(r\), the same as the column rank.

\section{Trace is contraction, and cyclicity is a legal index relabelling}

For endomorphisms \(A,B\),
\[
\begin{aligned}
  \operatorname{tr}(AB)
  &=(AB)_i{}^i\\
  &=A_i{}^jB_j{}^i\\
  &=B_j{}^iA_i{}^j\\
  &=\operatorname{tr}(BA).
\end{aligned}
\]
The middle equality only renames dummy indices and commutes scalar entries.  Consequently,
\[
  \operatorname{tr}(P^{-1}AP)
  =\operatorname{tr}(APP^{-1})
  =\operatorname{tr}A.
\]

For a rank-one operator
\[
  A=uv^*,
\]
one has
\[
  \boxed{
  \operatorname{tr}(uv^*)=v^*u.}
\]
Indeed,
\[
  (uv^*)_{ij}=u_i\overline{v_j},
\]
so
\[
  \operatorname{tr}(uv^*)
  =\sum_i u_i\overline{v_i}
  =v^*u.
\]
If \(u\) is a unit vector, the orthogonal projector onto its span is \(P=uu^*\), and
\[
  \operatorname{tr}P=1.
\]
More generally, an orthogonal projector has eigenvalues \(0,1\), hence
\[
  \operatorname{tr}P=\operatorname{rank}P.
\]
Trace counts the dimension of a projected subspace because it sums the projector's surviving modes.

\section{The Hilbert--Schmidt inner product leads to Schatten norms}

Matrices themselves form an inner-product space with
\[
  \langle A,B\rangle_{HS}
  =\operatorname{tr}(A^*B).
\]
Its norm is the Frobenius norm:
\[
\begin{aligned}
  \|A\|_F^2
  &=\operatorname{tr}(A^*A)\\
  &=\sum_{i,j}|a_{ij}|^2.
\end{aligned}
\]
If the singular values of \(A\) are
\[
  \sigma_1(A),\ldots,\sigma_r(A),
\]
then the Schatten norms are
\[
  \|A\|_{S_p}
  =
  \left(
  \sum_i\sigma_i(A)^p
  \right)^{1/p},
  \qquad1\le p<\infty,
\]
and
\[
  \|A\|_{S_\infty}=\max_i\sigma_i(A).
\]
The three common cases are
\[
  \|A\|_{S_1}=\text{nuclear norm},
  \qquad
  \|A\|_{S_2}=\|A\|_F,
  \qquad
  \|A\|_{S_\infty}=\|A\|_{op}.
\]

For
\[
  A=
  \begin{pmatrix}
    0&3\\
    4&0
  \end{pmatrix},
\]
one has
\[
  A^*A=
  \begin{pmatrix}
    16&0\\
    0&9
  \end{pmatrix},
\]
so the singular values are \(4,3\).  Therefore
\[
  \|A\|_{S_1}=7,
  \qquad
  \|A\|_{S_2}=5,
  \qquad
  \|A\|_{S_\infty}=4.
\]
Rank counts how many singular values are nonzero; Schatten norms measure their size with different emphases.

The Hilbert--Schmidt product also gives the estimate
\[
  |\operatorname{tr}(A^*B)|
  \le\|A\|_F\|B\|_F,
\]
which is ordinary Cauchy--Schwarz applied in the matrix space.

\section{Categorical trace closes an output wire into an input wire}

Let \(V\) be finite dimensional with basis \(e_i\) and dual basis \(e^i\).  The coevaluation map is
\[
  \operatorname{coev}:k\to V\otimes V^*,
  \qquad
  1\longmapsto\sum_i e_i\otimes e^i.
\]
The evaluation map is
\[
  \operatorname{ev}:V^*\otimes V\to k,
  \qquad
  \varphi\otimes v\longmapsto\varphi(v).
\]
For an endomorphism \(A:V\to V\), compose
\[
  k
  \xrightarrow{\operatorname{coev}}
  V\otimes V^*
  \xrightarrow{A\otimes I}
  V\otimes V^*
  \xrightarrow{\tau}
  V^*\otimes V
  \xrightarrow{\operatorname{ev}}
  k,
\]
where \(\tau(v\otimes\varphi)=\varphi\otimes v\).  Acting on \(1\), the composite gives
\[
\begin{aligned}
  1
  &\longmapsto
  \sum_i e_i\otimes e^i\\
  &\longmapsto
  \sum_i Ae_i\otimes e^i\\
  &\longmapsto
  \sum_i e^i\otimes Ae_i\\
  &\longmapsto
  \sum_i e^i(Ae_i).
\end{aligned}
\]
If
\[
  Ae_i=A_j{}^i e_j,
\]
then
\[
  \sum_i e^i(Ae_i)
  =\sum_iA_i{}^i
  =\operatorname{tr}A.
\]

The abstract operation of closing an output back into an input reproduces the ordinary matrix trace exactly.  Basis invariance is built into the evaluation--coevaluation pairing, while the coordinate calculation reveals the same contraction \(A_i{}^i\) encountered at the beginning of the chapter.